\documentclass[a4paper, 12pt]{amsart}

\usepackage[margin = 0.75in]{geometry}
\usepackage{imakeidx}
\makeindex

\newcommand{\defterm}[1]{\textbf{\textit{#1}}\index{#1}}

\newtheorem{theorem}{Theorem}[section]
\newtheorem{lemma}[theorem]{Lemma}
\newtheorem{proposition}[theorem]{Proposition}
\newtheorem{corollary}[theorem]{Corollary}

\theoremstyle{remark}
\newtheorem{remark}[theorem]{Remark}
\newtheorem{example}[theorem]{Example}

\begin{document}

\title{The Hilbert series of the focused arc algebra}
\author{Diego Salazar}
\date{\today}
\maketitle

An \defterm{$\mathbb{N}$-graded ring} is a ring $R$ together with a family $(R_n)_{n \in \mathbb{N}}$ of additive subgroups of $R$ such that $R = \bigoplus_{n \in \mathbb{N}} R_n$ and $R_mR_n \subseteq R_{m + n}$.
Let $k$ be a field.
A \defterm{$k$-algebra} is a commutative unital ring $A$ together with a ring homomorphism $k \hookrightarrow A$.
An \defterm{$\mathbb{N}$-graded $k$-algebra} is a $k$-algebra $A = \bigoplus_{n \in \mathbb{N}}A_n$ such that $A$ is an $\mathbb{N}$-graded ring, and each additive subgroup $A_n$ is a subspace of $A$.
If $a \in A_n$, we say $a$ has \defterm{weight} $n$.
This automatically implies $1 \in A_0$.

Let $A$ be an $\mathbb{N}$-graded $k$-algebra.
The \defterm{Hilbert-Poincaré series of $A$} is defined by
\begin{equation*}
  HP_A(t) = \sum_{n \in \mathbb{N}}\dim_k(A_n)t^n.
\end{equation*}
If each $A_n$ is finite dimensional, then $HP_A(t) \in \mathbb{N}[[t]]$.

A \defterm{differential $k$-algebra} is a $k$-algebra $A$ together with a $k$-linear derivation $\partial\in Der_k(A)$.
An \defterm{$\mathbb{N}$-graded differential $k$-algebra} is a differential $k$-algebra $A$ that is also an $\mathbb{N}$-graded $k$-algebra satisfying $\partial(A_n) \subseteq A_{n + 1}$.

Let $k$ be a field and let $A$ be a $k$-algebra.
Let $JA$ be the jet algebra of $A$.
We have a natural monomorphism $A \hookrightarrow JA$, so we can consider $A$ as a subset of $JA$ with $(A)_\partial = JA$, where the subscript $\partial$ indicates the differential ideal generated by the given subset.
We define
\begin{equation*}
  (JA)_n = span_k\{(\partial^{n_1}a_1)\dots(\partial^{n_s}a_s) \mid s, n_1, \dots, n_s \in \mathbb{N}, a_1, \dots, a_s \in A, n_1 + \dots + n_s = n\}.
\end{equation*}
Then
\begin{equation*}
  JA = \bigoplus_{n \in \mathbb{N}}(JA)_n,
\end{equation*}
and $JA$ becomes an $\mathbb{N}$-graded differential $k$-algebra.
We note that $(JA)_0 = A$.

Let $\mathfrak{p} \in Spec(A)$.
Then $A_\mathfrak{p}$ is a local ring with maximal ideal $\mathfrak{p}A_\mathfrak{p}$, and the quotient $\kappa(\mathfrak{p}) = A_\mathfrak{p}/\mathfrak{p}A_\mathfrak{p}$ is the \defterm{residue field of $A$ at $\mathfrak{p}$}.
We have a natural homomorphism $A \to \kappa(\mathfrak{p})$ which makes $\kappa(\mathfrak{p})$ into an $A$-algebra.
We can also view $JA$ as an $A$-algebra. 
The \defterm{focused arc algebra of $A$ at $\mathfrak{p}$} is the $\mathbb{N}$-graded $\kappa(\mathfrak{p})$-algebra
\begin{equation}
  \label{eq:1}
  J^\mathfrak{p}A = JA \otimes_A \kappa(\mathfrak{p})
\end{equation}
with operation given by
\begin{equation*}
  (a\otimes x)(b\otimes y) = (ab)\otimes(xy)
\end{equation*}
and with $\mathbb{N}$-grading given by
\begin{equation*}
(J^\mathfrak{p}A)_n = (JA)_n \otimes_A \kappa(\mathfrak{p}).
\end{equation*}
We note that $\dim_{\kappa(\mathfrak{p})}((J^\mathfrak{p}A)_0) = 1$ because $(J^\mathfrak{p}A)_0 = A \otimes_A \kappa(\mathfrak{p}) = \kappa(\mathfrak{p})$, and $J^\mathfrak{p}A$ does not inherit a derivation from $JA$.

We are interested in the case where
\begin{equation}
  \label{eq:2}
  A = k[X_1, \dots, X_n]/(F_1, \dots, F_m),
\end{equation}
where $F_1, \dots, F_m \in k[X_1, \dots, X_n]$ satisfy
\begin{equation}
  \label{eq:3}
  0 = F_r(0, \dots, 0), \qquad r = 1, \dots, m,
\end{equation}
and $\mathfrak{p} = (X_1 + (F_1, \dots, F_m), \dots, X_n + (F_1, \dots, F_m))$ is a maximal ideal.
In this case, the residue field is $\kappa(\mathfrak{p}) = k$, and the natural homomorphism $A \to k$ is just evaluation at $(0, \dots, 0)$.
Informally, we can think of $JA$ as being the following set of solutions
\begin{equation*}
  JA = \{x(t) \in k[[t]]^n \mid 0 = F_r(x(t)), r = 1, \dots, m \}
\end{equation*}
in the same way we can think of $A$ as
\begin{equation*}
  A = \{x \in k^n \mid 0 = F_r(x), r = 1, \dots, m\}.
\end{equation*}
Of course, this is just an intuition because the $k$-algebra $A$ may be nonreduced.

More precisely, we consider the polynomial $k$-algebra $k[X^{(i)}_j \mid i \in \mathbb{N}, j = 1, \dots, n]$ in countably many variables $X^{(i)}_j$ with derivation $\partial \in Der_k(k[X^{(i)}_j \mid i \in \mathbb{N}, j = 1, \dots, n])$ given by
\begin{equation}
  \label{eq:4}
  \partial X^{(i)}_j = (i + 1)X^{(i + 1)}_j, \qquad i \in \mathbb{N}, j = 1, \dots, n,
\end{equation}
where $char(k) = 0$.

To motivate this definition, we define:
\begin{equation}
  \label{eq:5}
  X_j(t) = X_j^{(0)} + X_j^{(1)}t + \dots \in k[X^{(i)}_j \mid i \in \mathbb{N}, j = 1, \dots, n][[t]], \qquad j = 1, \dots, n, \\
\end{equation}
We have a natural derivation $\partial_t$.
We wish to have
\begin{equation}
  \label{eq:6}
  \partial(F(X_1(t), \dots, X_n(t))) = \partial_t(F(X_1(t), \dots, X_n(t))), \qquad F \in k[X_1, \dots, X_n].
\end{equation}
Taking $F = X_j, j = 1, \dots, n$, we get \eqref{eq:4}.
With this definition, we get \eqref{eq:6} for $F \in k[X_1, \dots, X_n]$.

Identifying $X_j = X^{(0)}_j, j = 1, \dots, n$, we have
\begin{equation*}
  (X_1, \dots, X_n)_\partial = k[X^{(i)}_j \mid i \in \mathbb{N}, j = 1, \dots, n].
\end{equation*}

Following our intuition, we write
\begin{equation}
  \label{eq:7}
  F_r(X_1(t), \dots, X_n(t)) = F_r^{(0)} + F_r^{(1)}t + \dots, \qquad r = 1, \dots, m,
\end{equation}
where $F_r^{(l)} \in k[X^{(i)}_j \mid i \in \mathbb{N}, j = 1, \dots, n], l \in \mathbb{N}, r = 1 \dots, m$.
Then $(F^{(l)}_r \mid l \in \mathbb{N}, r = 1 \dots, m)$ is a differential ideal because $\partial(F_r(X(t))) = \partial_t(F_r(X(t))), r = 1, \dots, m$ implies
\begin{equation*}
  \partial(F_r^{(l)}) = (l + 1)F_r^{(l + 1)}, \qquad l \in \mathbb{N}, r = 1, \dots, m.
\end{equation*}
Thus
\begin{equation*}
  F_r^{(l)} = \frac{\partial^lF_r^{(0)}}{l!}, \qquad l \in \mathbb{N}, r = 1, \dots, m.
\end{equation*}
From \eqref{eq:7}, we see that $F_r^{(0)} = F_r, r = 1, \dots, m$.
We have proved that
\begin{equation*}
  (F^{(l)}_r \mid l \in \mathbb{N}, r = 1 \dots, m) = (F_1, \dots, F_m)_\partial.
\end{equation*}
We set
\begin{equation}
  \label{eq:8}
  JA = \frac{k[X^{(i)}_j \mid i \in \mathbb{N}, j = 1, \dots, n]}{(F^{(l)}_r \mid l \in \mathbb{N}, r = 1 \dots, m)} = \frac{k[X^{(i)}_j \mid i \in \mathbb{N}, j = 1, \dots, n]}{(F_1, \dots, F_m)_\partial},
\end{equation}
and we have a natural inclusion $A \hookrightarrow JA$ of $k$-algebras satisfying our desired universal property.

Setting
\begin{equation*}
  f^{(l)}_r = F_r^{(l)}|_{X^{(0)}_j = 0, j = 1 \dots, n} \in k[X^{(i)}_j \mid i \in \mathbb{Z}_+, j = 1, \dots, n], \qquad l \in \mathbb{Z}_+, r = 1 \dots, m,
\end{equation*}
we have
\begin{equation}
  \label{eq:9}
  J^0A = J^\mathfrak{p}A = JA \otimes_A k = \frac{k[X^{(i)}_j \mid i \in \mathbb{Z}_+, j = 1, \dots, n]}{(f^{(l)}_r \mid l \in \mathbb{Z}_+, r = 1, \dots, m)}.
\end{equation}

\begin{example}
  \label{exa:1}
  For $A = k[X]$, we have
  \begin{equation*}
    HP_{J^0A}(t) = \prod_{i \in \mathbb{Z}_+}\frac{1}{1 - t^i}.
  \end{equation*}
\end{example}

\begin{example}
  \label{exa:2}
  For $A = k[X]/(X^2)$, we have
  \begin{equation*}
    HP_{J^0A}(t) = \prod_{i \in \mathbb{Z}_+, i \equiv 1, 4 \bmod 5}\frac{1}{1 - t^i}.
  \end{equation*}
\end{example}

We could define $JA$ in another, more straightforward way following \eqref{eq:8} but modifying \eqref{eq:4}.
We again consider the polynomial $k$-algebra $k[X^{(i)}_j \mid i \in \mathbb{N}, j = 1, \dots, n]$ in countably many variables $X^{(i)}_j$ with derivation given by
\begin{equation}
  \label{eq:10}
  \partial X^{(i)}_j = X^{(i + 1)}_j, \qquad i \in \mathbb{N}, j = 1, \dots, n.
\end{equation}
Then we set
\begin{equation*}
  JA = \frac{k[X^{(i)}_j \mid i \in \mathbb{N}, j = 1, \dots, n]}{(F_1, \dots, F_m)_\partial}.
\end{equation*}
Again we have a natural inclusion $A \hookrightarrow JA$ of $k$-algebras satisfying our desired universal property.
We note we haven't assumed $char(k) = 0$ in this construction of the jet algebra.
If $char(k) = 0$, these two differential algebras are obviously isomorphic with isomorphism given by $X_j^{(i)} \mapsto i!X^{(i)}_j, i \in \mathbb{N}, j = 1, \dots, n$.

Alternatively, the graded structure on $JA$ can be defined more intrinsically via the natural action of the extension fields $K$ of $k$ on the arc parameter $t$.
Here, we again assume $char(k)$ might not be $0$.
Let $k \hookrightarrow K$ be a field extension of $k$.
For any $\lambda \in K^\times$, we wish to define a $K$-algebra automorphism $\Phi_\lambda \in Aut_K(K[X^{(i)}_j \mid i \in \mathbb{N}, j = 1, \dots, n])$ by setting its action on the generators:
\begin{equation}
  \label{eq:11}
  \Phi_\lambda(X^{(i)}_j) = \lambda^iX^{(i)}_j.
\end{equation}
We wish to have
\begin{equation}
  \label{eq:12}
  \Phi_\lambda(F(X_1(t), \dots, X_n(t))) = F(X_1(\lambda t), \dots, X_n(\lambda t)), \qquad F \in K[X_1, \dots, X_n].
\end{equation}
Indeed, taking $F = X_j, j = 1, \dots, n$ perfectly realizes this because
\begin{equation*}
  X_j(\lambda t) = \sum_{i = 0}^\infty (X^{(i)}_j\lambda^i)t^i,
\end{equation*}
and with this definition, we get \eqref{eq:12} for $F \in K[X_1, \dots, X_n]$.
Moreover, $G \in k[X^{(i)}_j \mid i \in \mathbb{N}, j = 1, \dots, n]$ has weight $d$ if and only if $\Phi_\lambda(G\otimes1) = G\otimes\lambda^d$ for all field extensions $K$ of $k$ and all $\lambda \in K^\times$.
It is necessary to consider field extensions $k \hookrightarrow K$ because if, for example, $k = \mathbb{Z}/2\mathbb{Z}$, then $\lambda^d = \lambda$ for all $d \in \mathbb{N}$ and $\lambda \in k^\times$.

To ensure $\Phi_\lambda$ is well-defined on the quotient $JA$, we must check that it preserves the defining ideal.
For the base polynomials $F_r \in k[X_1, \dots, X_n]$, the variables $X_j = X^{(0)}_j$ have weight $0$, so $\Phi_\lambda(F_r) = F_r$.
The derivation $\partial$ interacts with this automorphism via the commutation relation $\Phi_\lambda \circ \partial = \lambda \partial \circ \Phi_\lambda$.
This relation fundamentally holds for any derivation $\partial$ that maps $X^{(i)}_j$ to a scalar multiple of $X^{(i+1)}_j$, including the alternative definition \eqref{eq:10}.
For now, using the definition of $\partial$ from \eqref{eq:4}, we can prove this explicitly on the generators:
\begin{align*}
  \Phi_\lambda(\partial X^{(i)}_j) &= \Phi_\lambda((i+1)X^{(i+1)}_j) = (i+1)\lambda^{i+1}X^{(i+1)}_j, \\
  \lambda \partial(\Phi_\lambda(X^{(i)}_j)) &= \lambda \partial(\lambda^i X^{(i)}_j) = \lambda \cdot \lambda^i (i+1)X^{(i+1)}_j = (i+1)\lambda^{i+1}X^{(i+1)}_j.
\end{align*}
Because $\Phi_\lambda$ is a $k$-algebra automorphism and $\partial$ is a $k$-linear derivation, this relation extends from the generators to the entire polynomial ring.
Applying this relation recursively to $F^{(l)}_r = \partial^l F_r$ yields
\begin{equation*}
  \Phi_\lambda(F^{(l)}_r\otimes1) = F^{(l)}_r\otimes\lambda^l.
\end{equation*}
Since the generators of the differential ideal are simply scaled by $\lambda^l$, the ideal $(F_1, \dots, F_m)_\partial\otimes K$ is mapped entirely into itself.
Thus, $\Phi_\lambda$ descends to a well-defined automorphism $\Phi_\lambda: JA \otimes_k K \to JA \otimes_k K$. 
Therefore, an element $f \in JA$ is homogeneous of weight $d$ if and only if $\Phi_\lambda(f\otimes1) = f\otimes\lambda^d$ for all field extensions $k \hookrightarrow K$ and all $\lambda \in K^\times$.

Let $X$ be a $k$-scheme (not necessarily of finite type), i.e., we have a scheme morphism $X \to Spec(k)$, and let $k \hookrightarrow K$ be an extension of $k$.
The set of $K$-valued points is
\begin{equation*}
  X(K) = Hom_k(Spec(K), X) = \{(x, i) \mid x \in X, i: k(x) \to K\}.
\end{equation*}
Given $f \in \Gamma(X, \mathcal{O}_X)$, we define
\begin{align*}
  f_K: X(K) &\to K, \\
  f_K(x, i) &= i(f(x)),
\end{align*}
where
\begin{align*}
  \Gamma(X, \mathcal{O}_X) \to &\mathcal{O}_{X, x} \to k(x), \\
  f \mapsto &f_x \mapsto f(x).
\end{align*}

Let $X$ be a $k$-scheme of finite type.
Then $JX$ is a $k$-scheme (not generally of finite type) satisfying
\begin{equation*}
  JX(K) = Hom(Spec(K), JX) = Hom(Spec(K[[t]]), X)
\end{equation*}
for any extension $k \to K$ of $k$.

Given $\lambda \in K$, we have an endomorphism of $k$-algebras
\begin{align*}
  \varphi_\lambda: K[[t]] &\to K[[t]], \\
  \varphi_\lambda(t) &= \lambda t.
\end{align*}
We define an action
\begin{align*}
  K \times Hom_k(Spec(K[[t]]), X) &\to Hom_k(Spec(K[[t]]), X), \\
  (\lambda, \gamma) &\mapsto \lambda\cdot\gamma = \gamma\circ Spec(\varphi_\lambda).
\end{align*}

Finally, $f \in \Gamma(JX, \mathcal{O}_{JX})$ is homogeneous of weight $d$ if
\begin{equation*}
  f_K(\lambda\cdot(x, i)) = \lambda^df_K(x, i) \in K,
\end{equation*}
for all extensions $k \to K$ of $k$, $\lambda \in K$ and $(x, i) \in JX(K)$.

Maybe a better way of writing that is the following: if $x(t) \in Hom_k(Spec(K[[t]]), X)$
\begin{equation*}
  f_K(x(\lambda t)) = \lambda^df_K(x(t)).
\end{equation*}

CONJECTURE: $A = k[X]/F(X)$, where $F(X) = X^k + a_{k - 1}X^{k - 1} + \dots + a_iX^i$, $i \ge 1$.
\begin{equation*}
  J^0A = \text{identity of Gordon-Ramanujan-Poincare with k, i} 
\end{equation*}

Again, we consider $k$ a field.
Here, we again assume $char(k)$ might not be $0$.
For any $k$-algebra $K$ and any scalar $\lambda \in K^\times$, we have an automorphism of $K$-algebras
\begin{align*}
  \varphi_\lambda:K[[t]] &\to K[[t]], \\
  t &\mapsto \lambda t.
\end{align*}
This automorphism induces a natural scheme morphism over $K$
\begin{equation*}
  Spec(\varphi_\lambda): Spec(K[[t]]) \to Spec(K[[t]]),
\end{equation*}
which yields a natural bijection on the $K$-valued points:
\begin{equation*}
  \circ Spec(\varphi_\lambda): Hom_k(Spec(K[[t]]), X) \to Hom_k(Spec(K[[t]]), X).
\end{equation*}
Since $Hom_k(Spec(K[[t]]), X) = Hom_k(Spec(K), JX) = JX(K)$, we see that the group of units $K^\times$ acts naturally on $JX(K)$.
Recall that the multiplicative group scheme $\mathbb{G}_m = \operatorname{Spec}(k[\lambda, \lambda^{-1}])$ represents the functor that assigns to any $k$-algebra $K$ its group of units $K^\times$.
Thus, substituting $K^\times = \mathbb{G}_m(K)$, this defines a natural action of $\mathbb{G}_m$ on the functor of points of $JX$:
\begin{equation*}
  \mathbb{G}_m(K) \times JX(K) \to JX(K).
\end{equation*}
By Yoneda's Lemma, this uniquely determines a $k$-scheme action
\begin{equation*}
  \mu: \mathbb{G}_m \times_k JX \to JX.
\end{equation*}

Applying the global sections functor $\Gamma(-, \mathcal{O}_-)$, a group scheme action by $\mathbb{G}_m$ corresponds exactly to a $\mathbb{Z}$-grading on the coordinate ring $\Gamma(JX, \mathcal{O}_{JX})$. 
Specifically, the action $\mu$ induces a universal coaction homomorphism:
\begin{equation*}
  \mu^*: \Gamma(JX, \mathcal{O}_{JX}) \to \Gamma(JX, \mathcal{O}_{JX}) \otimes_k k[\lambda, \lambda^{-1}].
\end{equation*}
We say that $f \in \Gamma(JX, \mathcal{O}_{JX})$ is homogeneous of weight $d$ if
\begin{equation*}
  \mu^*(f) = f \otimes \lambda^d.
\end{equation*}
Equivalently, retreating to the functor of points, evaluating at a specific field extension $k \hookrightarrow K$ and an element $\lambda \in K^\times$ corresponds to a ring homomorphism $k[\lambda, \lambda^{-1}] \to K$.
This pushes the universal coaction forward to an automorphism of the extended coordinate ring:
\begin{equation*}
  \Gamma(\Phi_\lambda): \Gamma(JX, \mathcal{O}_{JX}) \otimes_k K \to \Gamma(JX, \mathcal{O}_{JX}) \otimes_k K.
\end{equation*}
In this space, $f$ is homogeneous of weight $d$ if it acts as expected on $f \otimes 1$:
\begin{equation*}
  \Gamma(\Phi_\lambda)(f\otimes1) = f\otimes\lambda^d.
\end{equation*}
(By abuse of notation, we often just write this as $\Gamma(\Phi_\lambda)(f) = \lambda^df$).
This definition operates entirely over arbitrary $k$-algebras $K$, perfectly bypassing finite-field collapses and avoiding the need to assume the scheme is reduced.

\subsection*{Summary (The functor of points)}
\label{sec:summary-the-functor}

In algebraic geometry, a scheme is not just its points over $k$.
It is defined by its points over all possible $k$-algebras $k \to K$.
Let $I \subseteq k[X_1, \dots, X_n]$ be an ideal, and let $X = Spec(k[X_1, \dots, X_n]/I)$.

\subsubsection*{1. The algebraic perspective (arbitrary rings)}
This perspective holds even if the base $k$ is merely a commutative ring rather than a field.
For any commutative $k$-algebra $K$ (not necessarily a field), the $K$-valued points of $X$ are given by the set
\begin{equation*}
  X(K) = Hom_{k-alg}(k[X_1, \dots, X_n]/I, K) \cong \{x \in K^n \mid g(x) = 0 \text{ for all } g \in I\}.
\end{equation*}
Note that when $K$ is not a field, this purely algebraic definition cannot be reduced to the geometric pairs $(x, i)$ defined earlier in the notes, because $Spec(K)$ is generally not a single point.
For $f + I \in k[X_1, \dots, X_n]/I$ and $x \in X(K)$, the usual evaluation $f_K(x) \in K$ is well-defined (independent of the representative $f$).
\begin{enumerate}
\item Two elements $f + I, g + I \in k[X_1, \dots, X_n]/I$ are equal if and only if $f_K(x) = g_K(x)$ for all $x \in X(K)$ and all $k$-algebras $K$.
\item Assuming $I$ is a homogeneous ideal, $f + I$ is homogeneous of degree $d$ if and only if $f_K(\lambda x) = \lambda^df_K(x)$ for all $k$-algebras $K$, all $x \in X(K)$, and all $\lambda \in K$.
\end{enumerate}
These equivalences hold unconditionally because one can choose the test algebra to be $K = k[X_1, \dots, X_n]/I$ itself, and evaluate at the generic point $x = (\bar{X}_1, \dots, \bar{X}_n) \in X(K)$, turning the conditions into algebraic tautologies.

\subsubsection*{2. The geometric perspective (fields)}
When $K$ is restricted to being a field extension of $k$, the spectrum $Spec(K)$ is just a single point.
In this case, a morphism $Spec(K) \to X$ corresponds to picking a topological point $x \in X$ and an embedding of its residue field $i: k(x) \hookrightarrow K$.
Thus, the $K$-valued points naturally identify with:
\begin{equation*}
  X(K) = \{(x, i) \mid x \in X, i: k(x) \to K\}.
\end{equation*}
This geometric definition provides clear intuition. In this setting, the evaluations $f_K(x)$ in properties (1) and (2) coincide exactly with the $f_K(x, i)$ defined earlier in the notes.
Furthermore, because we are studying reduced schemes over a field $k$, evaluating purely over field extensions $K$ is perfectly sufficient, meaning properties (1) and (2) still unconditionally hold in this restricted setting without needing to test over arbitrary rings.

\subsubsection*{Subtleties of evaluating over fields}
When restricting our test algebras $K$ to be fields, two important mathematical caveats arise regarding the equality of functions:

\begin{enumerate}
\item \textbf{Why ``reduced'' is required (The Nilpotent Blindspot):}
  Fields have no zero-divisors.
  Therefore, if a coordinate ring contains a non-zero nilpotent element $h \neq 0$ (e.g., $h^2 = 0$), evaluating $h$ at any point over any field $K$ forces $h(x)^2 = 0$, meaning $h(x) = 0 \in K$.
  Fields are mathematically blind to nilpotents.
  If a scheme is not reduced, two functions $f$ and $g$ differing by a nilpotent ($f - g = h$) will evaluate identically across all field extensions, even though $f \neq g$ algebraically.
  Thus, evaluating purely over fields only detects equality \textit{up to nilpotents}.
  This is precisely why we must be careful when defining homogeneity for the jet algebra $JA$, since quotienting by the differential ideal $(F_1, \dots, F_m)_\partial$ often introduces nilpotents.
  If we had defined homogeneity by evaluating functions at points via $f_K(\lambda x) = \lambda^d f_K(x)$ into field extensions $K$, the definition would catastrophically fail on non-reduced rings, falsely assigning arbitrary weights to nilpotents.

  However, our definition of homogeneity using $\Phi_\lambda$ completely bypasses this blindspot!
  By demanding $\Phi_\lambda(f) = \lambda^d f$ for all $\lambda \in K^\times$, we are evaluating the action of the ring automorphism $\Phi_\lambda: JA \otimes_k K \to JA \otimes_k K$.
  Because this algebraic identity lives entirely \textit{inside the extended coordinate ring} and never maps out to the scalar field $K$, non-zero nilpotents are never crushed to zero.
  Thus, even if $JA$ is non-reduced, taking field extensions $K$ is perfectly sufficient to separate the graded components.

For example, consider $X = \operatorname{Spec}(k[T]/(T^2))$.
The coordinate ring is not reduced because it contains the non-zero nilpotent element $T$.
For any field extension $K$, a point $x \in X(K)$ must satisfy $t^2 = 0$, which forces $t = 0$.
If we take the polynomials $f(T) = T$ and $g(T) = 0$, they evaluate identically to $0$ on all field-valued points.
The evaluation test over fields falsely concludes $f = g$, even though $T \neq 0$ in the ring.

\item \textbf{Why $\operatorname{char}(k) = 0$ is not enough and algebraic closure is needed:}
One might wonder if taking $K = k$ is sufficient when $\operatorname{char}(k) = 0$.
It is not.
Consider the affine scheme $X = \operatorname{Spec}(\mathbb{R}[T]/(T^2+1))$.
Here $k = \mathbb{R}$ has characteristic 0, and the coordinate ring is a field (hence perfectly reduced).
However, the set of $\mathbb{R}$-valued points is $X(\mathbb{R}) = \{t \in \mathbb{R} \mid t^2 + 1 = 0\} = \emptyset$.
Because the space of real points is empty, any two polynomials (e.g., $f(T) = T$ and $g(T) = 0$) vacuously evaluate to the exact same values on $X(\mathbb{R})$.
Yet $T \neq 0$ algebraically.
To guarantee that evaluating purely at $K=k$ points separates functions, $k$ must be \textbf{algebraically closed} (e.g., $k = \mathbb{C}$).
In this case, Hilbert's Nullstellensatz guarantees that $X(k)$ has enough points to detect equality of functions on reduced schemes.
\end{enumerate}

\printindex

\end{document}
